% Generated mechanically from frozen input inputs/p02/ja/tex/brauer.tex.
% Do not edit this generated file; edit the bound locale source instead.

% OK, start here.
%

\setcounter{chapter}{10}
\chapter{ブラウアー群}



\phantomsection
\label{brauer-section-phantom}


\section{序論}
\label{brauer-section-introduction}

\noindent
参考文献として、Cartan セミネールにおける Serre の講義
\cite{Serre-Cartan} がある。Serre はさらに
\cite{Deuring} と \cite{ANT} を参照している。本章ではいくつかの証明を
変更した。特に、Wedderburn の定理の証明には Rieffel の巧妙な議論を
用いた。この変更が改良になっているとは考えにくく、読者には Serre の
原典を読むことを強く勧める。


\section{非可換代数}
\label{brauer-section-algebras}

\noindent
$k$ を体とする。本章で $k$ 上の {\it 代数} $A$ とは、可換とは限らない
環 $A$ と環準同型 $k \to A$ の組であって、$k$ の像が $A$ の中心に
含まれ、かつ $1$ が $A$ の単位元に写るものをいう。{\it $A$-加群} とは、
$A$ の単位元が恒等写像として作用する右 $A$-加群をいう。

\begin{definition}
\label{brauer-definition-finite}
$A$ を $k$-代数とする。$\dim_k(A) < \infty$ のとき、$A$ は
{\it 有限} であるという。この場合、$[A : k] = \dim_k(A)$ と書く。
\end{definition}

\begin{definition}
\label{brauer-definition-skew-field}
{\it 斜体} とは、単位元 $1$ をもち、$1 \not = 0$ であり、すべての非零元が
乗法逆元をもつ、可換とは限らない環をいう。
\end{definition}

\noindent
斜体は、ある $k$ に対する $k$-代数である（例えば、その斜体に含まれる
素体を取ればよい）。以下では、斜体上の任意の加群が自由であることを
用いる。実際、ベクトルの極大線形独立集合は基底をなし、その存在は
Zorn の補題による。

\begin{definition}
\label{brauer-definition-simple}
$A$ を $k$-代数とする。$A$-加群 $M$ が非零で、その $A$-部分加群が
$0$ と $M$ のみであるとき、$M$ は {\it 単純} であるという。
$A$ の両側イデアルが $0$ と $A$ のみであるとき、$A$ は
{\it 単純} であるという。
\end{definition}

\begin{definition}
\label{brauer-definition-central}
$k$-代数 $A$ の中心が $k \to A$ の像に等しいとき、$A$ は
{\it 中心的} であるという。
\end{definition}

\begin{definition}
\label{brauer-definition-opposite}
$k$-代数 $A$ に対し、$A$ における積の順序を逆にして得られる
$k$-代数を $A^{op}$ と書く。これを {\it 反対代数} という。
\end{definition}




\section{Wedderburn の定理}
\label{brauer-section-wedderburn}

\noindent
次の巧妙な議論は Rieffel の論文 \cite{Rieffel} に見られる。
証明はこれ以上ないほど簡単である（Carl Faith の書評からの引用）。

\begin{lemma}
\label{brauer-lemma-rieffel}
$A$ を、$1$ をもち非自明な両側イデアルをもたない、可換とは限らない環と
する。$M$ を $A$ の非零右イデアルとし、右 $A$-加群とみなす。このとき
$A$ は $M$ の二重可換子と一致する。
\end{lemma}

\begin{proof}
$A' = \text{End}_A(M)$ と置けば、$M$ は左 $A'$-加群である。
$A'' = \text{End}_{A'}(M)$（$M$ の二重可換子）と置く。
$M$ を右 $A''$-加群とみなす\footnote{$a'' \in A''$ と $m \in M$ に
対して積 $m a'' \in M$ が定まるという意味である。特に、$A''$ の元を
左から作用する自己準同型として書いたときに得られる積とは、$A''$ の積の
順序が逆になる。}。
$mR(a) = ma$ で定まる自然な準同型 $R : A \to A''$ を考える。
$R(1) = \text{id}_M$ であり、$A$ は非自明な両側イデアルをもたないので、
$R$ は単射である。$R(M)$ が $A''$ の右イデアルであることを示す。
実際、$a'' \in A''$ とし、$m$ が $M$ に属するとき
$R(m)a'' = R(ma'')$ である。なぜなら、$M$ の任意の元 $n$ による
$M$ の {\it 左} 乗法は $A'$ の元を表し、したがって任意の $n$ が
$M$ に属するとき
$(nm)a'' = n(ma'')$ となるからである。
最後に、積イデアル $AM$ は両側イデアルであるから $A = AM$ である。
よって $R(A) = R(A)R(M)$ であり、$R(A)$ は $A''$ の右イデアルである。
しかも $R(A)$ は $A''$ の単位元を含むので、$R(A) = A''$ である。
\end{proof}

\begin{lemma}
\label{brauer-lemma-simple-module}
$A$ を $k$-代数とする。$A$ が有限ならば、次が成り立つ。
\begin{enumerate}
\item $A$ は単純加群をもつ。
\item 任意の非零加群は単純部分加群を含む。
\item $A$ 上の単純加群は $k$ 上有限次元である。
\item $M$ が単純 $A$-加群ならば、$\text{End}_A(M)$ は斜体である。
\end{enumerate}
\end{lemma}

\begin{proof}
$A$ 自身が非零 $A$-加群なので、(1) はもちろん (2) から従う。
(2) については、$k$-ベクトル空間として最小の（有限）次元をもつ
部分加群を取れば、それは単純である。巡回部分加群は有限次元なので、
そのような有限次元部分加群は存在する。$M$ が単純ならば
$mA \subset M$ は部分加群であり、これから (3) が従う。
$\text{End}_A(M)$ の任意の非零元は同型なので、(4) が成り立つ。
\end{proof}

\begin{theorem}
\label{brauer-theorem-wedderburn}
\begin{slogan}
体上の有限単純代数は、斜体上の行列代数である。
\end{slogan}
$A$ を単純有限 $k$-代数とする。このとき $A$ は、斜体である有限
$k$-代数 $K$ 上の行列代数である。
\end{theorem}

\begin{proof}
単純部分加群 $M \subset A$ を選べる。このとき
$k$-代数 $K = \text{End}_A(M)$ は斜体である。補題
\ref{brauer-lemma-simple-module} を参照されたい。
補題 \ref{brauer-lemma-rieffel} により
$A = \text{End}_K(M)$ である。$K$ は斜体であり、
$M$ は有限生成（$\dim_k(M) < \infty$ だから）なので、
$M$ は左 $K$-加群として有限自由である。したがって直ちに
$A \cong \text{Mat}(n \times n, K^{op})$ を得る。
\end{proof}






\section{代数に関する補題}
\label{brauer-section-lemmas}

\noindent
$A$ を $k$-代数とし、$B \subset A$ を部分代数とする。
{\it $A$ における $B$ の中心化代数} とは、部分代数
$$
C  = \{y \in A \mid xy = yx \text{ for all }x \in B\}.
$$
のことである。これは $k$-代数である。

\begin{lemma}
\label{brauer-lemma-centralizer}
$A$, $A'$ を $k$-代数とする。$B \subset A$, $B' \subset A'$ を
部分代数とし、それぞれの中心化代数を $C$, $C'$ とする。このとき
$A \otimes_k A'$ における $B \otimes_k B'$ の中心化代数は
$C \otimes_k C'$ である。
\end{lemma}

\begin{proof}
$C'' \subset A \otimes_k A'$ を $B \otimes_k B'$ の中心化代数とする。
$C \otimes_k C' \subset C''$ は明らかである。
逆に、$C''$ の任意の元は $B \otimes 1$ と可換するので
$C \otimes_k A'$ に属する。同様に $C'' \subset A \otimes_k C'$。
したがって
$C'' \subset C \otimes_k A' \cap A \otimes_k C' = C \otimes_k C'$。
\end{proof}

\begin{lemma}
\label{brauer-lemma-center-csa}
$A$ を有限単純 $k$-代数とする。このときその中心 $k'$ は
$k$ の有限体拡大である。
\end{lemma}

\begin{proof}
定理 \ref{brauer-theorem-wedderburn} により、$k$ 上有限なある斜体 $K$ に対して
$A = \text{Mat}(n \times n, K)$ と書ける。
補題 \ref{brauer-lemma-centralizer} により、$A$ の中心は
$k \otimes_k k'$ である。ここで $k' \subset K$ は $K$ の中心である。
斜体の中心は体なので、結論を得る。
\end{proof}

\begin{lemma}
\label{brauer-lemma-generate-two-sided-sub}
$V$ を $k$-ベクトル空間とする。$K$ を斜体である中心的
$k$-代数とする。$W \subset V \otimes_k K$ を両側
$K$-部分ベクトル空間とする。このとき $W$ は、左 $K$-ベクトル空間として
$W \cap (V \otimes 1)$ によって生成される。
\end{lemma}

\begin{proof}
$v \otimes 1 \in W$ を満たす $v \in V$ によって生成される
$k$-部分ベクトル空間を $V' \subset V$ とする。このとき
$V' \otimes_k K \subset W$ であり、
$$
W/(V' \otimes_k K)  \subset  (V/V') \otimes_k K.
$$
$\overline{v} \in V/V'$ が非零ベクトルで、
$\overline{v} \otimes 1$ が $W/(V' \otimes_k K)$ に含まれるならば、
$\overline{v}$ の持ち上げ $v \in V$ に対して $v \otimes 1 \in W$ と
なる。これは $V'$ の構成に反する。したがって $V$ を $V/V'$ で、
$W$ を $W/(V' \otimes_k K)$ で置き換えてよく、$W$ が非零ならば
$W \cap (V \otimes 1)$ も非零であることを示せば十分である。

\medskip\noindent
これを見るため、非零元 $w \in W$ を
$w = \sum_{i = 1, \ldots, n} v_i \otimes k_i$ と書けるもののうち
$n$ が最小になるように取る。右から $k_1^{-1}$ を掛けて
$k_1 = 1$ と仮定してよい。$n = 1$ なら
$v_1 \otimes 1 \in W$ なので結論を得る。
$n > 1$ ならば、任意の $c \in K$ に対して
$$
c w - w c = \sum\nolimits_{i = 2, \ldots, n} v_i \otimes (c k_i - k_i c) \in W
$$
であり、$n$ の最小性から $c k_i - k_i c = 0$ である。
したがって $k_i$ は $K$ の中心、すなわち仮定により $k$ に属する。
ゆえに $w = (v_1 + \sum k_i v_i) \otimes 1$ となり、
$n$ の最小性に反する。
\end{proof}

\begin{lemma}
\label{brauer-lemma-generate-two-sided-ideal}
$A$ を $k$-代数とする。$K$ を斜体である中心的 $k$-代数とする。
このとき任意の両側イデアル $I \subset A \otimes_k K$ は、$A$ の
ある両側イデアル $J \subset A$ に対して $J \otimes_k K$ の形である。
特に $A$ が単純ならば $A \otimes_k K$ も単純である。
\end{lemma}

\begin{proof}
$J = \{a \in A \mid a \otimes 1 \in I\}$ と置く。これは $A$ の
両側イデアルである。補題 \ref{brauer-lemma-generate-two-sided-sub} により
$I = J \otimes_k K$ である。
\end{proof}

\begin{lemma}
\label{brauer-lemma-matrix-algebras}
$R$ を可換とは限らない環とし、$n \geq 1$ を整数とする。
$R_n = \text{Mat}(n \times n, R)$ と置く。
\begin{enumerate}
\item 関手 $M \mapsto M^{\oplus n}$ と
$N \mapsto Ne_{11}$ は、圏の互いに準逆な同値
$\text{Mod}_R \leftrightarrow \text{Mod}_{R_n}$ を定める。
\item $R_n$ の任意の両側イデアルは、$R$ のある両側イデアル $I$ に
対して $IR_n$ の形である。
\item $R_n$ の中心は $R$ の中心に等しい。
\end{enumerate}
\end{lemma}

\begin{proof}
(1) は明らかである。$J \subset R_n$ が両側イデアルならば
$J = \bigoplus e_{ii}Je_{jj}$ であり、すべての直和因子
$e_{ii}Je_{jj}$ は互いに等しく、$R$ の両側イデアル $I$ になる。
これで (2) が従う。(3) は明らかである。
\end{proof}

\begin{lemma}
\label{brauer-lemma-simple-module-unique}
$A$ を有限単純 $k$-代数とする。
\begin{enumerate}
\item 同型を除いてただ一つの単純 $A$-加群 $M$ が存在する。
\item 任意の有限 $A$-加群は単純加群のコピーの直和である。
\item 二つの有限 $A$-加群が同型であることと、それらの $k$ 上の
次元が等しいことは同値である。
\item $A = \text{Mat}(n \times n, K)$ で、$K$ が $k$ の有限斜体拡大
ならば、$M = K^{\oplus n}$ は単純 $A$-加群であり、
$\text{End}_A(M) = K^{op}$ である。
\item $M$ が単純 $A$-加群ならば、$L = \text{End}_A(M)$ は $k$ 上
有限な斜体で、$M$ に左から作用する。また
$A = \text{End}_L(M)$ であり、$A$ と $L$ の中心は一致する。
さらに $[A : k] [L : k] = \dim_k(M)^2$ である。
\item 有限 $A$-加群 $N$ に対し、代数 $B = \text{End}_A(N)$ は
(5) の斜体 $L$ 上の行列代数である。さらに
$\text{End}_B(N) = A$ である。
\end{enumerate}
\end{lemma}

\begin{proof}
定理 \ref{brauer-theorem-wedderburn} により、$k$ のある有限斜体拡大 $K$ に
対して $A = \text{Mat}(n \times n, K)$ と書ける。
補題 \ref{brauer-lemma-matrix-algebras} により、$A$-加群の圏は
$K$-加群の圏と同値である。$K$ 上の任意の加群は自由なので、
(1)、(2)、(3) が従う。(4) は、この同値が $K$-加群 $K$ を
$M = K^{\oplus n}$ に移すことから従う。(5) で
$M = K^{\oplus n}$ を用いれば $L = K^{op}$ を得る。
$L = K^{op}$ の中心に関する主張は補題
\ref{brauer-lemma-matrix-algebras} から従う。
$\text{End}_L(M)$ に関する主張は $M$ の具体的な形から従い、
次元公式も明らかである。(6) は、$N$ が単純加群のコピーの直和に
同型であることから従う。
\end{proof}

\begin{lemma}
\label{brauer-lemma-tensor-simple}
$A$, $A'$ を二つの単純 $k$-代数とし、その一方が $k$ 上有限かつ
中心的であるとする。このとき $A \otimes_k A'$ は単純である。
\end{lemma}

\begin{proof}
$A'$ が $k$ 上有限かつ中心的であるとする。
定理 \ref{brauer-theorem-wedderburn} により
$A' = \text{Mat}(n \times n, K')$ と書く。
$K'$ の中心は $k$ なので、補題
\ref{brauer-lemma-generate-two-sided-ideal} により
$A \otimes_k K'$ は単純である。したがって補題
\ref{brauer-lemma-matrix-algebras} により
$A \otimes_k A' = \text{Mat}(n \times n, A \otimes_k K')$ は単純である。
\end{proof}

\begin{lemma}
\label{brauer-lemma-tensor-central-simple}
$k$ 上の有限中心単純代数のテンソル積は、有限、中心的かつ単純である。
\end{lemma}

\begin{proof}
補題 \ref{brauer-lemma-centralizer} と \ref{brauer-lemma-tensor-simple} を合わせればよい。
\end{proof}

\begin{lemma}
\label{brauer-lemma-base-change}
$A$ を有限中心単純 $k$-代数とし、$k'/k$ を体拡大とする。このとき
$A' = A \otimes_k k'$ は有限中心単純 $k'$-代数である。
\end{lemma}

\begin{proof}
補題 \ref{brauer-lemma-centralizer} と \ref{brauer-lemma-tensor-simple} を合わせればよい。
\end{proof}

\begin{lemma}
\label{brauer-lemma-inverse}
$A$ を有限中心単純 $k$-代数とする。このとき
$n = [A : k]$ とすれば
$A \otimes_k A^{op} \cong \text{Mat}(n \times n, k)$ である。
\end{lemma}

\begin{proof}
補題 \ref{brauer-lemma-tensor-central-simple} により
$A \otimes_k A^{op}$ は単純である。したがって写像
$$
A \otimes_k A^{op} \longrightarrow \text{End}_k(A),\quad
a \otimes a' \longmapsto (x \mapsto axa')
$$
は単射である。矢印の両辺の次元が等しいので、結論を得る。
\end{proof}





\section{体のブラウアー群}
\label{brauer-section-brauer}

\noindent
$k$ を体とする。$k$ 上の二つの有限中心単純代数 $A$ と $B$ を考える。
ある $n, m > 0$ が存在して
$\text{Mat}(n \times n, A) \cong \text{Mat}(m \times m, B)$
が $k$-代数として成り立つとき、$A$ と $B$ は {\it 相似} であるという。

\begin{lemma}
\label{brauer-lemma-similar}
相似関係について、次が成り立つ。
\begin{enumerate}
\item 相似関係は、$k$ 上の有限中心単純代数の同型類の集合上の
同値関係を定める。
\item 各相似類は、同型を除いて一意な $k$ の有限中心斜体拡大を含む。
\item $K$, $K'$ を $k$ 上の有限中心斜体とし、
$A = \text{Mat}(n \times n, K)$ および
$B = \text{Mat}(m \times m, K')$ とする。このとき $A$ と $B$ が
相似であることと、$K \cong K'$ が $k$-代数として成り立つことは同値である。
\end{enumerate}
\end{lemma}

\begin{proof}
Wedderburn の定理（定理 \ref{brauer-theorem-wedderburn}）により、有限中心単純代数は
常に有限中心斜体上の行列代数として書ける。したがって第三の主張を
示せば十分である。そのためには、
$A = \text{Mat}(n \times n, K) \cong
\text{Mat}(m \times m, K') = B$ ならば $K \cong K'$ であることを
示せばよい。$A$ の単純加群 $M$ に対して、補題
\ref{brauer-lemma-simple-module-unique} により
$\text{End}_A(M) = K^{op}$ である。よって
$A \cong B$ から $K^{op} \cong (K')^{op}$ が従い、結論を得る。
\end{proof}

\noindent
$k$ 上の二つの有限中心単純代数 $A$, $B$ に対し、補題
\ref{brauer-lemma-tensor-central-simple} によればテンソル積
$A \otimes_k B$ も有限中心単純代数である。
さらに $A$ が $A'$ と相似ならば、テンソル積と行列代数を取る操作は
可換するので、$A \otimes_k B$ は $A' \otimes_k B$ と相似である。
したがってテンソル積は有限中心単純代数の同値類上に、明らかに結合的かつ
可換な演算を定める。最後に補題 \ref{brauer-lemma-inverse} により
$A \otimes_k A^{op}$ は行列代数と同型、すなわち
$A \otimes_k A^{op}$ は $k$ の相似類に属する。こうして可換群を得る。

\begin{definition}
\label{brauer-definition-brauer-group}
$k$ を体とする。上で定めた有限中心単純 $k$-代数の相似類からなる
可換群を、$k$ の {\it ブラウアー群} という。記号は
$\text{Br}(k)$ である。
\end{definition}

\noindent
任意の体の準同型 $k \to k'$ に対し、補題
\ref{brauer-lemma-base-change} により群準同型
$$
\text{Br}(k) \longrightarrow \text{Br}(k'),\quad
A \longmapsto A \otimes_k k'
$$
を得る。言い換えれば、$\text{Br}(-)$ は体の圏から可換群の圏への
関手である。体のブラウアー群が零であることと、任意の有限中心斜体拡大
$k \subset K$ が自明であることは同値である。

\begin{lemma}
\label{brauer-lemma-brauer-algebraically-closed}
代数閉体のブラウアー群は零である。
\end{lemma}

\begin{proof}
$k \subset K$ を有限中心斜体拡大とする。任意の元 $x \in K$ に対し、
部分環 $k[x] \subset K$ は可換で有限かつ $k$ 上整な部分代数なので体である。
Algebra, Lemma \ref{algebra-lemma-integral-over-field} を参照されたい。
$k$ は代数閉なので $k[x] = k$ である。$x$ は任意だったから $k = K$。
\end{proof}

\begin{lemma}
\label{brauer-lemma-dimension-square}
$A$ を体 $k$ 上の有限中心単純代数とする。このとき $[A : k]$ は
平方数である。
\end{lemma}

\begin{proof}
補題 \ref{brauer-lemma-brauer-algebraically-closed} により
$A \otimes_k \overline{k}$ は $\overline{k}$ 上の行列代数なので、主張が従う。
\end{proof}




\section{Skolem--Noether}
\label{brauer-section-skolem-noether}



\begin{theorem}
\label{brauer-theorem-skolem-noether}
$A$ を有限中心単純 $k$-代数とし、$B$ を単純 $k$-代数とする。
$f, g : B \to A$ を二つの $k$-代数準同型とする。このとき、任意の
$b \in B$ に対して $f(b) = xg(b)x^{-1}$ を満たす可逆元
$x \in A$ が存在する。
\end{theorem}

\begin{proof}
単純 $A$-加群 $M$ を選び、$L = \text{End}_A(M)$ と置く。
補題 \ref{brauer-lemma-simple-module} と \ref{brauer-lemma-simple-module-unique} により、
$L$ は中心 $k$ をもつ斜体で、$M$ に左から作用する。
$M$ 上に二つの $B \otimes_k L^{op}$-加群構造を
$m \cdot_1 (b \otimes l) = lmf(b)$ と
$m \cdot_2 (b \otimes l) = lmg(b)$ により定める。
補題 \ref{brauer-lemma-tensor-simple} により
$k$-代数 $B \otimes_k L^{op}$ は単純である。$B$ は単純であり
$k$-代数準同型 $B \to A$ が存在するので、$B$ は有限である。
したがって $B \otimes_k L^{op}$ は有限単純であり、補題
\ref{brauer-lemma-simple-module-unique} により $M$ 上の二つの
$B \otimes_k L^{op}$-加群構造は同型である。
よってこれらの作用を絡み合わせる $\varphi : M \to M$ が存在する。
特に $\varphi$ は $L$ の可換子に属するので、補題
\ref{brauer-lemma-simple-module-unique} により、$\varphi$ は
ある $x \in A$ による乗法である。
定義を書き下せば、$x$ が求める解であることが分かる。
\end{proof}

\begin{lemma}
\label{brauer-lemma-automorphism-inner}
$A$ を有限中心単純 $k$-代数とする。$A$ の任意の自己同型は内部自己同型である。
特に $\text{Mat}(n \times n, k)$ の任意の自己同型は内部自己同型である。
\end{lemma}

\begin{proof}
Skolem--Noether の定理（定理 \ref{brauer-theorem-skolem-noether}）を適用すればよい。
\end{proof}



\section{中心化定理}
\label{brauer-section-centralizer}


\begin{theorem}
\label{brauer-theorem-centralizer}
$A$ を $k$ 上の有限中心単純代数とし、$B$ を $A$ の単純部分代数とする。
このとき次が成り立つ。
\begin{enumerate}
\item $A$ における $B$ の中心化代数 $C$ は単純である。
\item $[A : k] = [B : k][C : k]$ である。
\item $A$ における $C$ の中心化代数は $B$ である。
\end{enumerate}
\end{theorem}

\begin{proof}
この証明では補題 \ref{brauer-lemma-simple-module-unique} の結果を自由に用いる。
単純 $A$-加群 $M$ を選び、$L = \text{End}_A(M)$ と置く。
すると $L$ は中心 $k$ をもつ斜体で $M$ に左から作用し、
$A = \text{End}_L(M)$ である。
このとき $M$ は右 $B \otimes_k L^{op}$-加群であり、
$C = \text{End}_{B \otimes_k L^{op}}(M)$ である。
補題 \ref{brauer-lemma-tensor-simple} により代数
$B \otimes_k L^{op}$ は単純なので、再び補題
\ref{brauer-lemma-simple-module-unique} により $C$ は単純である。

\medskip\noindent
$k$ 上有限なある斜体 $K$ に対して
$B \otimes_k L^{op} = \text{Mat}(m \times m, K)$ と書く。
$M$ が単純 $B \otimes_k L^{op}$-加群 $K^{\oplus m}$ の $n$ 個の
コピーの直和に同型ならば、同じ補題により
$C = \text{Mat}(n \times n, K^{op})$ である。したがって
$\dim_k(M) = nm [K : k]$、
$[B : k] [L : k] = m^2 [K : k]$、
$[C : k] = n^2 [K : k]$、および同じ補題により
$[A : k] [L : k] = \dim_k(M)^2$ である。これから (2) が従う。

\medskip\noindent
(3) を示す。$A$ における $C$ の中心化代数を $C'$ とする。
$C \subset A$ に (2) を適用すれば $[B : k] = [C' : k]$ となる。
一方（明らかに $B \subset C')$ なので、(3) が従う。
\end{proof}

\begin{lemma}
\label{brauer-lemma-when-tensor-is-equal}
$A$ を有限中心単純 $k$-代数とし、$B$ を $A$ の単純部分代数とする。
$B$ が中心的 $k$-代数ならば、$C$ を $A$ における $B$ の
（中心単純な）中心化代数として $A = B \otimes_k C$ である。
\end{lemma}

\begin{proof}
定理 \ref{brauer-theorem-centralizer} により
$\dim_k(A) = \dim_k(B \otimes_k C)$ である。
補題 \ref{brauer-lemma-tensor-simple} によりテンソル積は単純である。
したがって自然な写像 $B \otimes_k C \to A$ は単射であり、
次元が等しいので同型である。
\end{proof}

\begin{lemma}
\label{brauer-lemma-self-centralizing-subfield}
$A$ を有限中心単純 $k$-代数とする。$K \subset A$ が部分体ならば、
次の条件は同値である。
\begin{enumerate}
\item $[A : k] = [K : k]^2$。
\item $K$ は自身の中心化代数である。
\item $K$ は極大可換部分環である。
\end{enumerate}
\end{lemma}

\begin{proof}
定理 \ref{brauer-theorem-centralizer} により (1) と (2) は同値である。
(3) と (2) の同値性は明らかである。
\end{proof}

\begin{lemma}
\label{brauer-lemma-maximal-subfield}
\begin{slogan}
有限中心斜体の次元は、その任意の極大部分体の次元の平方である。
\end{slogan}
$A$ を $k$ 上の有限中心斜体とする。このとき任意の極大部分体
$K \subset A$ は $[A : k] = [K : k]^2$ を満たす。
\end{lemma}

\begin{proof}
補題 \ref{brauer-lemma-self-centralizing-subfield} の特別な場合である。
\end{proof}





\section{分解体}
\label{brauer-section-splitting}


\begin{definition}
\label{brauer-definition-splitting}
$A$ を有限中心単純 $k$-代数とする。$A \otimes_k k'$ が $k'$ 上の
行列代数であるとき、体拡大 $k'/k$ は $A$ を {\it 分解する}、
または $k'$ は $A$ の {\it 分解体} であるという。
\end{definition}

\noindent
これは、$A$ の類が写像
$\text{Br}(k) \to \text{Br}(k')$ の下で零に写ると言い換えてもよい。

\begin{theorem}
\label{brauer-theorem-splitting}
$A$ を有限中心単純 $k$-代数とし、$k'/k$ を有限体拡大とする。
次の条件は同値である。
\begin{enumerate}
\item $k'$ は $A$ を分解する。
\item $A$ と相似な有限中心単純代数 $B$ で、
$k' \subset B$ かつ $[B : k] = [k' : k]^2$ を満たすものが存在する。
\end{enumerate}
\end{theorem}

\begin{proof}
(2) を仮定する。$B \otimes_k k'$ が行列代数であることを示せば十分である。
$B \otimes_k B^{op} \cong \text{End}_k(B)$ である。
補題 \ref{brauer-lemma-self-centralizing-subfield} により、
$k'$ は $B^{op}$ における $k'$ の中心化代数なので、
$B \otimes_k k'$ は
$B \otimes_k B^{op} = \text{End}_k(B)$ における
$k \otimes k'$ の中心化代数である。この中心化代数は、
埋め込み $k' \to B$ によって $B$ を $k'$-ベクトル空間とみたときの
$\text{End}_{k'}(B)$ にほかならない。よって結論を得る。

\medskip\noindent
(1) を仮定する。これは、ある $k'$-ベクトル空間 $V$ に対して
$A \otimes_k k' \cong \text{End}_{k'}(V)$ という同型があることを意味する。
$B$ を $\text{End}_k(V)$ における $A$ の可換子とする。
$k'$ は $B$ に含まれることに注意する。
補題 \ref{brauer-lemma-when-tensor-is-equal} により、$A$ と $B$ の類の和は
$\text{Br}(k)$ で零である。定理 \ref{brauer-theorem-centralizer} の次元公式から
$$
[B : k] [A : k] =
\dim_k(V)^2 =
[k' : k]^2 \dim_{k'}(V)^2 =
[k' : k]^2 [A : k].
$$
を得る。したがって $[B : k] = [k' : k]^2$ である。
以上で $A$ のブラウアー類の逆元について結果を示した。しかし $k'$ が
$A$ のブラウアー類を分解することと、反対代数のブラウアー類を分解する
ことは同値なので、いずれにせよ結論を得る。
\end{proof}

\begin{lemma}
\label{brauer-lemma-maximal-subfield-splits}
$k$ 上の有限中心斜体 $K$ の極大部分体は、$K$ の分解体である。
\end{lemma}

\begin{proof}
補題 \ref{brauer-lemma-maximal-subfield} と定理
\ref{brauer-theorem-splitting} を合わせればよい。
\end{proof}

\begin{lemma}
\label{brauer-lemma-splitting-field-degree}
$k$ 上の有限中心斜体 $K$ を考え、$d^2 = [K : k]$ とする。
$K$ の任意の有限分解体 $k'$ に対し、次数 $[k' : k]$ は $d$ で割り切れる。
\end{lemma}

\begin{proof}
定理 \ref{brauer-theorem-splitting} により、$K$ のブラウアー類に属し
$[B : k] = [k' : k]^2$ を満たす有限中心単純代数 $B$ が存在する。
補題 \ref{brauer-lemma-similar} により、ある $n$ に対して
$B = \text{Mat}(n \times n, K)$ である。
したがって $[k' : k]^2 = n^2d^2$ であり、結論を得る。
\end{proof}

\begin{proposition}
\label{brauer-proposition-separable-splitting-field}
$k$ 上の有限中心斜体 $K$ を考える。$k$ 上分離的な極大部分体
$k \subset k' \subset K$ が存在する。特に、任意のブラウアー類は
有限分離分解体をもつ。
% P02-E0001
\end{proposition}

\begin{proof}
任意のブラウアー類は $k$ 上の有限中心斜体によって表されるので、
第二の主張は補題 \ref{brauer-lemma-maximal-subfield-splits} により第一の主張から
従う。

\medskip\noindent
第一の主張を示すため、分離的な部分体 $k' \subset K$ が与えられたとする。
$K$ における $k'$ の中心化代数 $K'$ は中心 $k'$ をもち、問題は
$k'$ 上分離的な $K'$ の極大部分体を見つけることに帰着する。
したがって $k \not = K$ のとき、$k$ 上分離的で
$x \not \in k$ を満たす元 $x \in K$ を見つければ十分である。
標数零ではこれは明らかである。よって $k$ の標数は $p > 0$ と仮定してよい。
基礎体 $k$ が有限ならば、有限体の拡大は常に分離的なので、この場合も
明らかである。したがって $k$ は正標数の無限体と仮定してよい。

\medskip\noindent
背理法のため、$K$ のどの元も $k$ 上分離的でないと仮定する。
% P02-E0002
Fields, Section \ref{fields-section-algebraic} の議論により、
これは任意の $x \in K$ の最小多項式が
$T^q - a$ の形であることを意味する。ここで $q$ は $p$ の冪であり
$a \in k$ である。$K$ の各元の最小多項式の次数が
$\leq \dim_k(K)$ であることは明らかなので、すべての $x \in K$ に対して
$x^q \in k$ となる固定した $p$-冪 $q$ が存在する。

\medskip\noindent
写像
$$
(-)^q : K \longrightarrow K
$$
を考え、$a_1 = 1$ を満たす $K$ の $k$-基底
$\{a_1, \ldots, a_n\}$ を用いて書き下す。すると
$$
(\sum x_i a_i)^q = \sum f_i(x_1, \ldots, x_n)a_i.
$$
$K$ 上の乗法は $k$-双線形なので、各 $f_i$ は
$x_1, \ldots, x_n$ の多項式である（詳細は省略する）。
上で選んだ $q$ と $k$ が無限であることから、$i \geq 2$ に対して
$f_i$ は恒等的に零である。したがって $k$ をその代数閉包
$\overline{k}$ に拡大しても零のままである。しかし
$K \otimes_k \overline{k}$ は行列代数である（例えば補題
\ref{brauer-lemma-base-change} と \ref{brauer-lemma-brauer-algebraically-closed} による）。
したがって $q$ 乗が中心に属さない元（例えば $e_{11}$）が存在する。
これは求める矛盾である。
\end{proof}

\noindent
以上の結果により、有限中心単純代数を次のように特徴づけられる。

\begin{lemma}
\label{brauer-lemma-finite-central-simple-algebra}
$k$ を体とする。$k$-代数 $A$ に対し、次の条件は同値である。
\begin{enumerate}
\item $A$ は有限中心単純 $k$-代数である。
\item $A$ は有限次元 $k$-ベクトル空間で、$k$ は $A$ の中心であり、
$A$ は非自明な両側イデアルをもたない。
\item ある $d \geq 1$ が存在して
$A \otimes_k \bar k \cong \text{Mat}(d \times d, \bar k)$ である。
\item ある $d \geq 1$ が存在して
$A \otimes_k k^{sep} \cong \text{Mat}(d \times d, k^{sep})$ である。
\item ある $d \geq 1$ と有限 Galois 拡大 $k'/k$ が存在して
$A \otimes_k k' \cong \text{Mat}(d \times d, k')$ である。
\item ある $n \geq 1$ と $k$ 上の有限中心斜体 $K$ が存在して
$A \cong \text{Mat}(n \times n, K)$ である。
\end{enumerate}
整数 $d$ を $A$ の {\it 次数} という。
\end{lemma}

\begin{proof}
(1) と (2) の同値性は定義から従う。第
\ref{brauer-section-algebras} 節を参照されたい。
(1) を仮定する。命題
\ref{brauer-proposition-separable-splitting-field} により、$A$ には分離分解体
$k \subset k'$ が存在する。もちろん $k'/k$ の Galois 閉包も分解体である。
したがって (1) は (5) を含意する。(5) $\Rightarrow$ (4)
$\Rightarrow$ (3) は明らかである。(3) を仮定する。このとき
$A \otimes_k \overline{k}$ は、例えば補題
\ref{brauer-lemma-matrix-algebras} により有限中心単純
$\overline{k}$-代数である。これは直ちに $A$ が有限中心単純
$k$-代数であることを含意する。最後に、(1) と (6) の同値性は
Wedderburn の定理、すなわち定理 \ref{brauer-theorem-wedderburn} である。
\end{proof}
